\documentclass[11pt]{article}
\usepackage{amsmath,amssymb,amsthm}
\usepackage[margin=1.1in]{geometry}
\usepackage{booktabs}

\newtheorem{proposition}{Proposition}
\newtheorem{remark}{Remark}

\title{Poisson Regression, and the GLM Unification}
\author{Yuval Lavie}
\date{\today}

\begin{document}
\maketitle

\begin{abstract}
The conditional-distribution recipe of the logistic note, applied to
count data: $Y|X{=}x \sim \mathrm{Poisson}(\lambda(x))$ with log-linear
rate. The per-sample gradient is once again \emph{residual times
features} --- and this third instance is no coincidence: for any
exponential-family response with its canonical link, the gradient of the
conditional NLL is $(\hat\mu - y)\,x$. That unification --- generalized
linear models --- is derived here, making the linear, logistic, and
Poisson notes three rows of one table.
\end{abstract}

\section{The model}

Counts $y \in \{0,1,2,\dots\}$ with
\begin{equation}
  \Pr(Y=y \mid X=x)
  = e^{-\lambda(x)} \frac{\lambda(x)^{y}}{y!},
  \qquad
  \lambda(x) = e^{w^\top x} > 0 ,
  \label{eq:model}
\end{equation}
the exponential (log link) keeping the rate positive exactly as the
sigmoid kept a probability in $(0,1)$. Poisson facts:
$\mathbb{E}[Y|x] = \mathrm{Var}(Y|x) = \lambda(x)$ (mean equals
variance --- visible in the script's first panel), and the log-linear
form means each unit increase of a feature \emph{multiplies} the expected
count by $e^{w_j}$.

\section{Conditional MLE}

The per-sample NLL, dropping the $\log y_i!$ term (constant in $w$):
\begin{equation}
  \mathcal{L}_i(w)
  \;=\; \lambda_i - y_i \log\lambda_i
  \;=\; e^{w^\top x_i} - y_i\, w^\top x_i .
  \label{eq:nll}
\end{equation}
Differentiate: $\nabla_w \mathcal{L}_i = e^{w^\top x_i} x_i - y_i x_i$,
hence
\begin{equation}
  \nabla_w \mathcal{L}_i \;=\; \bigl(\lambda_i - y_i\bigr)\, x_i .
  \label{eq:grad}
\end{equation}
Residual times features, again. Stationarity
$\sum_i e^{w^\top x_i} x_i = \sum_i y_i x_i$ is transcendental --- no
closed form, as in the logistic case --- but the Hessian
\begin{equation}
  \nabla^2 \mathcal{L} \;=\; \frac1n X^\top
  \operatorname{diag}(\lambda_i)\, X \;\succeq\; 0
  \label{eq:hessian}
\end{equation}
is PSD with weights equal to the conditional \emph{variances} (for
Poisson, the rates themselves), so the loss is convex and Newton/IRLS
--- each step a weighted least-squares solve --- converges quadratically
(8 steps vs GD's thousands in the script).

\section{The GLM unification}

The pattern is now unmistakable:
\begin{center}
\begin{tabular}{llll}
\toprule
response $Y|x$ & link ($\theta = w^\top x$) & mean $\hat\mu(\theta)$ &
per-sample gradient\\
\midrule
Gaussian & identity & $\theta$ & $(\theta - y)\,x$\\
Bernoulli & logit & $\sigma(\theta)$ & $(\sigma(\theta) - y)\,x$\\
Poisson & log & $e^{\theta}$ & $(e^{\theta} - y)\,x$\\
\bottomrule
\end{tabular}
\end{center}

\begin{proposition}[Canonical-link GLM gradient]
Let the conditional distribution be an exponential family in natural
parameter $\theta$,
$\;p(y \mid \theta) = h(y)\, \exp\bigl(\theta\, y - A(\theta)\bigr)$,
with $\theta = w^\top x$ (the \emph{canonical link}). Then the per-sample
NLL is $\mathcal{L}_i(w) = A(w^\top x_i) - y_i\, w^\top x_i$ and
\begin{equation}
  \boxed{\;\nabla_w \mathcal{L}_i
  \;=\; \bigl(\hat\mu(w^\top x_i) - y_i\bigr)\, x_i ,
  \qquad \hat\mu = A' = \mathbb{E}[Y\mid\theta].\;}
  \label{eq:glm}
\end{equation}
\end{proposition}

\begin{proof}
$\nabla_w \mathcal{L}_i = A'(w^\top x_i)\,x_i - y_i x_i$, and the
exponential-family identity $A'(\theta) = \mathbb{E}[Y\mid\theta]$
(differentiate $\int h e^{\theta y - A} = 1$ in $\theta$) identifies
$A'$ as the mean function. Convexity is automatic:
$A''(\theta) = \mathrm{Var}(Y\mid\theta) \ge 0$, so the Hessian is
$\tfrac1n X^\top \operatorname{diag}(\mathrm{Var}_i) X \succeq 0$ ---
which is also why the IRLS weights are always the conditional variances.
\end{proof}

Check the table against the identity: Gaussian, $A = \theta^2/2$,
$A' = \theta$; Bernoulli, $A = \log(1+e^{\theta})$, $A' = \sigma$;
Poisson, $A = e^{\theta}$, $A' = e^{\theta}$. One theorem, all three
gradients --- and the reason ``prediction minus target, pushed along the
features'' keeps reappearing: it is not a coincidence of examples but
the derivative of the exponential family's log-normalizer. Any new
response distribution in this family (exponential, gamma, multinomial
--- see the softmax note) joins the table for free.

\section{Experiment}

$n = 5000$, $w^\ast = (0.3, 0.8, -0.5)$, $x \sim \mathcal{N}(0, I_2)$,
counts drawn from \eqref{eq:model} (seed 7). Newton/IRLS (8 steps) and GD
(4000 steps) agree on the optimum, recovering $w^\ast$ to sampling
accuracy; hand and autograd SGD trajectories coincide to $10^{-10}$; the
calibration panel plots fitted $e^{\hat w^\top x}$ against true
$\lambda(x)$ on fresh points --- the estimated object is the \emph{rate
function}, the Poisson sibling of logistic regression's $p(x)$.

\end{document}
